%=============================================================================
% Wave 4, Iteration 14: Materials Science Update
% AGENT: MATERIALS SCIENCE (W4i14) | Model: Opus 4.6
%
% INHERITED FROM W4i13 MatSci:
%   - Crystalline Si: Q>10^10 at 7K, 45x margin, 430x SNR improvement
%   - Nb3Sn Phase III: $75-230M savings at 4K
%   - O-doped Nb: backup for Phase I
%   - AlN PnC: built-in piezoelectricity for readout
%   - 4H-SiC: path to Q~10^10 at room temperature (speculative)
%   - KILLED: diamond NV, topological insulators, metamaterials, levitated nanoparticles
%   - Phase I Q0 sensitivity scaling question (1/Q0 vs 1/sqrt(Q0))
%
% W4i14 MANDATE:
%   1. Fabrication protocols for crystalline Si PnC membranes
%   2. Nb surface optimisation: O-doping vs N-doping quantitative comparison
%   3. 4H-SiC path to Q~10^10 --- CRITICAL UPDATE from Phys Rev Applied 2025
%   4. AlN piezo readout development and interface piezoelectricity
%   5. High-Tc superconductors for room-temp-adjacent SRF
%   6. Novel phononic band engineering
%   7. 2D material resonators (graphene, MoS2) --- KILL assessment
%   8. 3C-SiC electromechanical systems --- new competitor?
%
% Date: 20 March 2026
%=============================================================================

\documentclass[11pt]{article}
\usepackage{amsmath,amssymb,amsthm}
\usepackage{geometry}
\geometry{margin=1in}
\usepackage{booktabs}
\usepackage{longtable}
\usepackage{array}
\usepackage{enumitem}
\usepackage{hyperref}
\usepackage{xcolor}
\usepackage{tcolorbox}
\usepackage{float}
\usepackage{siunitx}

\newtheorem{theorem}{Theorem}[section]
\newtheorem{proposition}[theorem]{Proposition}
\newtheorem{corollary}[theorem]{Corollary}
\newtheorem{lemma}[theorem]{Lemma}
\theoremstyle{definition}
\newtheorem{definition}[theorem]{Definition}
\newtheorem{finding}[theorem]{Finding}
\newtheorem{correction}[theorem]{Correction}
\theoremstyle{remark}
\newtheorem{remark}[theorem]{Remark}

\newcommand{\ee}[1]{\times 10^{#1}}
\newcommand{\Meff}{M_{\rm eff}}
\newcommand{\Qpsi}{Q_\psi}
\newcommand{\dpsi}{\delta\psi_{\rm EM}}
\newcommand{\psigrav}{\psi_{\rm grav}}
\newcommand{\psiEM}{\delta\psi_{\rm EM}}
\newcommand{\lambdaone}{|\lambda-1|}
\newcommand{\lambdabound}{1.56\ee{-9}}
\newcommand{\Tc}{T_{\rm c}}
\newcommand{\lambdaL}{\lambda_{\rm L}}
\newcommand{\FOMpassive}{\mathrm{FOM}_{\rm passive}}
\newcommand{\FOMMEMS}{\mathrm{FOM}_{\rm MEMS}}
\newcommand{\lam}{|\lambda-1|}
\newcommand{\Qmech}{Q_{\rm mech}}
\newcommand{\SNR}{\mathrm{SNR}}
\newcommand{\confirmed}[1]{\textcolor{blue}{\textbf{CONFIRMED:} #1}}
\newcommand{\corrected}[1]{\textcolor{red}{\textbf{CORRECTED:} #1}}
\newcommand{\caution}[1]{\textcolor{orange}{\textbf{CAUTION:} #1}}
\newcommand{\honest}[1]{\textcolor{violet}{\textbf{HONEST ASSESSMENT:} #1}}
\newcommand{\killed}[1]{\textcolor{red!70!black}{\textbf{KILLED:} #1}}
\newcommand{\verdict}[1]{\textcolor{green!50!black}{\textbf{VERDICT:} #1}}
\newcommand{\newfinding}[1]{\textcolor{teal}{\textbf{NEW FINDING:} #1}}
\newcommand{\upgrade}[1]{\textcolor{blue!70!black}{\textbf{UPGRADE:} #1}}

\title{Wave 4 Iteration 14: Materials Science Update\\
\large Fabrication Protocols for Crystalline Si PnC,\\
Nb O-Doping vs N-Doping Quantitative Comparison,\\
4H-SiC Critical Reassessment,\\
Interface Piezoelectricity and 3C-SiC Electromechanics,\\
2D Resonator Kill Assessment, High-T$_c$ SRF Survey}
\author{DFD Research Programme --- Materials Science Agent, W4i14}
\date{20 March 2026}

\begin{document}
\maketitle
\tableofcontents
\newpage

%=============================================================================
\section{Executive Summary: Top 5 Findings}
\label{sec:top5}
%=============================================================================

\begin{tcolorbox}[colback=blue!5!white, colframe=blue!50!black,
  title=\textbf{W4i14 TOP 5 FINDINGS --- Ranked by Programme Impact}]
\begin{enumerate}

\item \textbf{FINDING 1 (CORRECTION): 4H-SiC soft-clamping path to
$Q \sim 10^{10}$ at room temperature is BLOCKED --- stress-free
monolithic 4H-SiC cannot exploit dissipation dilution.}
Sementilli et al.\ (Phys.\ Rev.\ Applied 24, 034005, September 2025)
demonstrate $Q_{\rm int}$ up to $2\ee{5}$ at room temperature in
monolithic 4H-SiC nanostrings electrochemically etched from
monocrystalline wafers. Crucially, these devices have
\textit{essentially zero built-in stress}. The authors explicitly
state that ``tensile-strain related strategies to enhance the
mechanical quality factor, including dissipation dilution, soft
clamping, and strain engineering, cannot be exploited in monolithic
4H-SiC resonators due to the absence of mechanical stress.''

\textbf{Derivation chain:} Dissipation dilution factor
$D_Q = \sigma L^2/(E t^2)$ requires tensile pre-stress $\sigma > 0$.
For monolithic 4H-SiC: $\sigma \approx 0$ (no epitaxial mismatch
in homoepitaxial etch), so $D_Q \approx 1$. The W4i13 projection
of $D_Q \sim 10^3$ applied to $Q_{\rm int} = 2\ee{7}$ giving
$Q \sim 2\ee{10}$ was \textbf{physically incorrect} for monolithic
4H-SiC. The $Q_{\rm int} = 2\ee{7}$ value from W4i13 (Sementilli
et al., Nano Letters 2025) is also now superseded: the Phys.\ Rev.\
Applied paper reports $Q_{\rm int}$ up to $2\ee{5}$ for nanostrings
approaching the thermoelastic limit at $f > 10$~MHz.

\corrected{W4i13 Finding 5 overstated the 4H-SiC path. The projection
of $Q \sim 10^{10}$ at room temperature via dissipation dilution is
\textbf{blocked} for monolithic 4H-SiC. A heterogeneous approach
(4H-SiC film on stressed Si or Si$_3$N$_4$ substrate) could restore
the dilution pathway but has not been demonstrated. 4H-SiC
\textbf{downgraded from WATCH to LONG-TERM SPECULATIVE}.}

However, the exceptionally low Akhiezer damping in 4H-SiC
(Gr\"uneisen parameter $\gamma \approx 0.35$, lowest of any
semiconductor) means that \textit{if} a stressed SiC film can be
fabricated (e.g., via heteroepitaxial growth on Si), the Akhiezer
$Q \cdot f$ limit would be $\sim 5\ee{14}$~Hz, comparable to
Si$_3$N$_4$. This remains a material science challenge, not a
physics impossibility.

\item \textbf{FINDING 2: Nb O-doping vs N-doping --- quantitative
head-to-head comparison establishes O-doping as a viable Phase~I
alternative with distinct advantages at high fields.}
Hu, Kim \& Bafia (Appl.\ Phys.\ Lett.\ 128, 112601, March 2026)
present the first systematic comparison of nitrogen and oxygen
interstitial doping in 1.3~GHz Nb SRF cavities under identical
processing conditions with material characterisation on cavity
cutouts.

Key quantitative results:
\begin{itemize}[nosep]
\item N-doping: increases average superconducting gap $\Delta$
  at $\sim 10\times$ lower concentration than O-doping,
  suggesting stronger per-atom effect (likely via H-trapping
  suppression of Nb nanohydrides).
\item N-doping: reduces quasiparticle density-of-states smearing
  vs electropolished baseline, lowering $R_{\rm BCS}$.
\item O-doping: achieves comparable $R_s$ reduction to N-doping
  but with improved $Q_0$ stability at higher accelerating
  gradients ($E_{\rm acc} > 20$~MV/m), where N-doped cavities
  show Q-slope onset.
\item O-doping: provides a distinct optimisation axis orthogonal
  to N-doping, enabling combinatorial surface engineering.
\end{itemize}

\textbf{DFD Phase~I impact:} For the SQMS Phase~I 4-cavity array,
the Q-ratio diagnostic (psi-prediction 0.49 vs BCS-prediction 3.41)
requires stable $Q_0$ across a frequency range. O-doped cavities'
improved high-field stability could reduce systematic uncertainty in
the Q-ratio measurement. \textbf{Recommendation: request SQMS provide
one O-doped and one N-doped cavity pair within the 4-cavity array
for direct comparison.} Cost increment: zero if within SQMS programme
scope (O-doping uses existing furnace infrastructure with ${\rm O}_2$
gas instead of ${\rm N}_2$).

\verdict{O-doped Nb upgraded from ``backup'' (W4i13) to
``co-primary'' for Phase~I. The N-doped/O-doped comparison within
the same 4-cavity array provides an additional systematic cross-check
on the Q-ratio diagnostic at no extra cost.}

\item \textbf{FINDING 3: Nb$_3$Sn co-sputtering from composite
targets achieves $T_c = 17.78$~K with improved surface homogeneity
--- Phase~III fabrication pathway strengthened.}
Fermilab (2025, fermilab-pub-25-0591-td) demonstrates Nb$_3$Sn
coating of a 2.6~GHz SRF cavity by DC cylindrical magnetron
co-sputtering from a composite Nb-Sn target, achieving:
\begin{itemize}[nosep]
\item $T_c = 17.78$~K (close to highest reported, cf.\ bulk
  $T_c = 18.3$~K)
\item Sn content 32--42 at.\% on beam tubes and equator
\item Optimal anneal: 600$^\circ$C/6h + 950$^\circ$C/1h
\item Light Sn recoating + re-anneal enhances RF performance
  via improved surface homogeneity
\end{itemize}

Separately, grain boundary (GB) research (Acta Materialia, 2020;
ongoing at JLab) establishes that Sn segregation at GBs is the
primary $Q_0$-limiting defect. Optimised processes avoiding GB Sn
segregation (confirmed by STEM-EDS) correlate with high cavity
performance. The co-sputtering approach provides better control over
Sn stoichiometry than vapor diffusion, potentially reducing GB
segregation.

\textbf{Derivation chain for Phase~III cost update:}
W4i13 estimated \$75--230M savings. With co-sputtering enabling
more reproducible coatings, the yield assumption improves from
$\sim 50\%$ (vapor diffusion) to $\sim 70\%$ (co-sputtering,
projected). At 70\% yield, 256 good cavities require
$256/0.7 = 366$ coating runs vs $256/0.5 = 512$ at 50\% yield.
At \$5K/coating: \$1.83M vs \$2.56M --- saving \$0.73M on
coating alone. The dominant saving remains cryogenics
(4~K cryocoolers vs 20~mK dilution refrigerators): \$13--26M vs
\$100--256M, giving \$75--230M saving (unchanged from W4i13).

\confirmed{Co-sputtering strengthens the Nb$_3$Sn Phase~III pathway
by improving surface homogeneity and potentially yield. The \$75--230M
cryogenic saving estimate from W4i13 is UNCHANGED but the fabrication
risk is reduced.}

\item \textbf{FINDING 4: Interface piezoelectricity at Al/Si junctions
discovered as a loss channel in superconducting devices --- direct
relevance to DFD MEMS-SRF integration.}
Jahanbani et al.\ (Nature Communications, December 2025) report the
first observation of interface piezoelectricity at aluminum-silicon
junctions in superconducting devices. The effective electromechanical
coupling is $K^2 \approx (3 \pm 0.4) \times 10^{-5}\%$, comparable
to weakly piezoelectric substrates.

\textbf{DFD programme implications:}
\begin{enumerate}[nosep]
\item \textbf{Loss channel for crystalline Si MEMS at mK:} If a
  crystalline Si PnC resonator is mounted in a cryostat with Al
  wire bonds or metallisation, the interface piezoelectric coupling
  can mediate energy loss from mechanical modes to spurious
  electromagnetic modes. The quality factor limitation scales as
  $Q_{\rm piezo}^{-1} \sim K^2 \cdot p_{\rm surface}$, where
  $p_{\rm surface}$ is the surface participation ratio. For a
  PnC defect mode with $p_{\rm surface} \sim 10^{-4}$ (soft-clamped):
  $Q_{\rm piezo} \sim 1/((3\ee{-7}) \times 10^{-4}) \sim 3\ee{10}$.
  This is \textit{comparable to the target $Q = 10^{10}$}, meaning
  interface piezoelectricity could be a limiting loss mechanism
  for DFD's cryogenic crystalline Si resonators.

\item \textbf{Mitigation}: Use phononic bandgap engineering to
  suppress acoustic mode density at the Si-metal interface. The
  soft-clamped PnC design already provides this for the defect
  mode, but spurious modes in the bandgap wings could couple.
  Alternative: use non-metallic (optical) readout to eliminate
  Al-Si interfaces entirely.

\item \textbf{Opportunity}: The same interface piezoelectricity
  provides a new transduction mechanism for MEMS readout ---
  direct electromechanical coupling at Si-Al interfaces without
  requiring an intrinsically piezoelectric material like AlN.
  $K^2 = 3\ee{-7}$ is weak but nonzero, enabling capacitive
  transduction schemes with superconducting microwave circuits.
\end{enumerate}

\caution{Interface piezoelectricity could limit crystalline Si
PnC resonator $Q$ at cryogenic temperatures to $\sim 3\ee{10}$
if Al metallisation is used. This does not affect the 45$\times$
$Q/\Qpsi$ margin (threshold is $2.2\ee{8}$, so even $Q = 3\ee{10}$
gives $136\times$ margin), but it may affect the 430$\times$ SNR
improvement claimed in W4i13. The thermal noise reduction (43$\times$
from $T$) is unaffected; only the $Q$ enhancement factor could be
capped at $30\times$ instead of $100\times$ if $Q$ saturates at
$3\ee{10}$ rather than exceeding $10^{10}$.}

\item \textbf{FINDING 5: 3C-SiC membrane crystals demonstrate 21
high-$Q$ multimode electromechanical coupling with superconducting
circuits --- a dark horse for DFD multimode psi-detection.}
Li et al.\ (Nature Communications 16, 1207, January 2025) and a
follow-up (npj Quantum Information, 2026) demonstrate:
\begin{itemize}[nosep]
\item 3C-SiC square membranes with non-uniform built-in tensile stress
\item 21 high-$Q$ mechanical modes controlled from a single membrane
\item Degeneracy-breaking via anisotropic stress in cubic crystal
\item Direct microwave electromechanical coupling to superconducting
  circuits
\item Extremely low pure dephasing rates
\item Group delays up to $\sim 1$~hour (indicating very high
  coherence)
\end{itemize}

Unlike monolithic 4H-SiC (Finding~1), 3C-SiC \textit{does} have
built-in tensile stress from heteroepitaxial growth on Si substrates
(lattice mismatch $\sim 20\%$ after relaxation), enabling dissipation
dilution. The multimode character is attractive for DFD because
multiple mechanical modes spanning a frequency range could be used
for broadband psi-field spectroscopy, potentially improving the
in-gap diagnostic by providing frequency-resolved $\Qpsi$ data.

\honest{3C-SiC membranes are a genuine new option for DFD MEMS
detection, combining tensile stress (unlike 4H-SiC), multimode
operation, and native electromechanical coupling. However, the
$Q$ values have not been characterised in the same way as Si$_3$N$_4$
or crystalline Si PnC resonators, and the demonstrated $Q$ per mode
is not yet at the $10^9$--$10^{10}$ level needed for DFD. Monitor
for $Q$ characterisation of individual defect modes in soft-clamped
3C-SiC PnC geometry.}

\verdict{3C-SiC added to materials watch list at MEDIUM priority
(above 4H-SiC, below crystalline Si). Unique multimode + electromech
coupling makes it a candidate for Phase~III broadband psi-spectroscopy.}

\end{enumerate}
\end{tcolorbox}


%=============================================================================
\section{Finding 1: 4H-SiC Path to $Q \sim 10^{10}$ --- BLOCKED}
\label{sec:sic_correction}
%=============================================================================

\subsection{W4i13 Claim vs New Evidence}

W4i13 Finding 5 projected that monolithic 4H-SiC could achieve
$Q \sim 10^{10}$ at room temperature by applying dissipation dilution
($D_Q \sim 10^3$) to the demonstrated $Q_{\rm int} = 2\ee{7}$. This
projection was based on analogy with Si$_3$N$_4$, where soft clamping
improved $Q$ from $\sim 10^5$ (bulk) to $\sim 10^9$ (soft-clamped PnC).

The critical error: Si$_3$N$_4$ has $\sigma \sim 1$~GPa built-in
tensile stress from LPCVD deposition. Monolithic 4H-SiC fabricated
by electrochemical etching from homoepitaxial wafers has
$\sigma \approx 0$ --- no lattice mismatch, no deposition-induced
stress.

\subsection{Quantitative Analysis}

The dissipation dilution factor for a doubly-clamped beam:
\begin{equation}
D_Q = \frac{\sigma L^2}{E t^2} \cdot \frac{1}{n^2 \pi^2},
\label{eq:dq_beam}
\end{equation}
where $\sigma$ is the tensile stress, $L$ is length, $E$ is Young's
modulus ($\sim 450$~GPa for 4H-SiC), $t$ is thickness, and $n$ is
the mode number. For $\sigma \to 0$: $D_Q \to 0$, meaning
\textit{no dilution is possible}.

The Phys.\ Rev.\ Applied paper confirms: the best $Q$ values
are $Q_{\rm int} \sim 2\ee{5}$ at room temperature approaching the
thermoelastic limit at $f > 10$~MHz, and the $Q \cdot f$ product
exceeds $10^{12}$~Hz but is far below the $Q \cdot f \sim 6\ee{14}$~Hz
achieved by soft-clamped Si$_3$N$_4$.

\subsection{Salvage Pathways for SiC}

Three potential routes to restore the $Q \sim 10^{10}$ target:

\begin{enumerate}[nosep]
\item \textbf{Heteroepitaxial 3C-SiC on Si}: 3C-SiC films grown on
  Si have $\sim 20\%$ lattice mismatch, producing built-in tensile
  stress after partial relaxation. This enables dissipation dilution,
  as demonstrated by Li et al.\ (Nat.\ Commun.\ 2025). However,
  3C-SiC has higher defect density than 4H-SiC, potentially limiting
  $Q_{\rm int}$.

\item \textbf{Externally stressed 4H-SiC}: Apply tensile stress via
  a frame structure (similar to how Si$_3$N$_4$ gets stress from
  the Si substrate). Would require a stressed 4H-SiC-on-insulator
  (SiCOI) platform, which is under development for quantum photonics
  but not yet optimised for nanomechanics.

\item \textbf{Akhiezer-limited bulk acoustic wave (BAW) modes}:
  Instead of flexural modes, use BAW modes in thick SiC resonators.
  The Akhiezer $Q \cdot f$ limit for 4H-SiC is $\sim 5\ee{14}$~Hz
  (from $\gamma = 0.35$). At $f = 1$~GHz: $Q_{\rm Akhiezer} \sim
  5\ee{5}$; at $f = 10$~MHz: $Q_{\rm Akhiezer} \sim 5\ee{7}$.
  Still far below $10^{10}$ without dilution.
\end{enumerate}

\begin{finding}[4H-SiC Correction]
\label{find:sic_correction}
\corrected{W4i13 Finding 5 was physically incorrect: monolithic 4H-SiC
has zero built-in stress and cannot benefit from dissipation dilution.
The projection of $Q \sim 10^{10}$ at room temperature is blocked.
Demonstrated $Q_{\rm int} = 2\ee{5}$ (Phys.\ Rev.\ Applied 2025),
not $2\ee{7}$ as cited in W4i13 (the $2\ee{7}$ figure appears to be
from a different fabrication route or measurement frequency regime).
4H-SiC downgraded from WATCH to LONG-TERM SPECULATIVE.}

\newfinding{3C-SiC (cubic polytype) retains a path to high $Q$ via
built-in stress from heteroepitaxial growth. 3C-SiC promoted to
WATCH (see Finding~5).}
\end{finding}


%=============================================================================
\section{Finding 2: Nb O-Doping vs N-Doping --- Quantitative Comparison}
\label{sec:o_doping}
%=============================================================================

\subsection{Head-to-Head Results (Hu, Kim \& Bafia, APL March 2026)}

This paper provides the first controlled comparison:

\begin{table}[H]
\centering
\caption{N-doping vs O-doping of 1.3~GHz Nb SRF cavities.}
\renewcommand{\arraystretch}{1.3}
\begin{tabular}{p{4cm}cc}
\toprule
\textbf{Property} & \textbf{N-doped} & \textbf{O-doped} \\
\midrule
Effective concentration for $\Delta$ increase & Lower ($\sim 10\times$) & Higher \\
Per-atom effect on $\Delta$ & Stronger & Weaker \\
$R_{\rm BCS}$ reduction & Significant & Comparable \\
$Q_0$ at low field & Excellent & Excellent \\
$Q_0$ stability at $E_{\rm acc} > 20$~MV/m & Q-slope onset & More stable \\
H-trapping (nanohydride suppression) & Primary mechanism & Secondary \\
Residual resistance & Low & Low \\
Processing infrastructure & ${\rm N}_2$ gas furnace & ${\rm O}_2$ gas furnace \\
\bottomrule
\end{tabular}
\end{table}

\subsection{Physical Mechanism}

N-doping works primarily by trapping interstitial hydrogen, preventing
formation of lossy niobium nanohydrides (${\rm NbH}_x$) in the
RF penetration depth ($\sim 40$~nm). The superconducting gap
$\Delta$ increases because the pair-breaking effect of nanohydrides
is eliminated. N is more effective per atom because its electronic
configuration allows stronger H-binding.

O-doping works via a related but distinct mechanism: oxygen
interstitials modify the electronic density of states near the Fermi
level, reducing quasiparticle generation at RF frequencies. The
improved high-field stability may arise from oxygen's ability to pin
magnetic flux vortices more effectively than nitrogen, reducing
flux-flow dissipation at high $E_{\rm acc}$.

\subsection{DFD Phase~I Strategy}

\begin{finding}[O-Doping Upgrade for Phase~I]
\label{find:o_doping}
\upgrade{O-doped Nb upgraded from ``backup'' to ``co-primary'' for
Phase~I.}

\textbf{Recommended 4-cavity Phase~I configuration:}
\begin{itemize}[nosep]
\item Cavity 1 (1.3~GHz): N-doped Nb (SQMS baseline recipe)
\item Cavity 2 (1.3~GHz): O-doped Nb (new APL recipe)
\item Cavity 3 (3.9~GHz): N-doped Nb
\item Cavity 4 (3.9~GHz): O-doped Nb
\end{itemize}

This configuration enables:
\begin{enumerate}[nosep]
\item The primary Q-ratio diagnostic at two frequencies
  (psi: 0.49 vs BCS: 3.41)
\item A doping-method cross-check: any psi-signal should appear
  identically in N-doped and O-doped pairs, while BCS systematics
  will differ
\item Direct measurement of the high-field Q-slope difference
  between N and O doping
\end{enumerate}

\textbf{Cost increment: \$0} if O-doped cavities are requested within
the SQMS programme scope. The furnace infrastructure is identical;
only the process gas changes.
\end{finding}


%=============================================================================
\section{Finding 3: Nb$_3$Sn Co-Sputtering Advancement}
\label{sec:nb3sn_cosputtering}
%=============================================================================

\subsection{Fermilab Co-Sputtering Results (2025)}

The Fermilab co-sputtering programme (fermilab-pub-25-0591-td)
represents a significant advancement over vapor diffusion:

\begin{table}[H]
\centering
\caption{Nb$_3$Sn fabrication methods comparison.}
\renewcommand{\arraystretch}{1.3}
\begin{tabular}{p{3cm}cc}
\toprule
\textbf{Property} & \textbf{Vapor Diffusion} & \textbf{Co-Sputtering} \\
\midrule
$T_c$ achieved & 18.0--18.3~K & 17.78~K \\
Sn composition control & Limited (gradient) & 32--42 at.\% (tunable) \\
Surface homogeneity & Moderate & Improved (post recoat) \\
GB Sn segregation & Common & Reducible \\
Anneal protocol & 1100$^\circ$C/3h & 600$^\circ$C/6h + 950$^\circ$C/1h \\
Scalability to 256 cavities & Challenging & More reproducible \\
Cu substrate compatible & No & Yes (recipe optimised) \\
\bottomrule
\end{tabular}
\end{table}

\subsection{Grain Boundary Engineering}

STEM-EDS analysis confirms that Sn segregation at grain boundaries
is the primary $Q_0$-limiting defect in Nb$_3$Sn coatings. Optimal
anneal protocols (as in the co-sputtering work) avoid GB segregation,
and cavities without GB segregation show minimal Q-factor degradation.

The co-sputtering approach allows independent control of Sn content
and anneal temperature, providing a two-dimensional optimisation space
(Sn composition $\times$ anneal protocol) for minimising GB defects.

\subsection{Compact Cryomodule Development}

RadiaBeam Technologies (in collaboration with Fermilab, 2025) has
developed a compact cryomodule for conduction-cooled Nb$_3$Sn cavities,
designed for a 4.5-cell 650~MHz cavity on a 4~K pulse-tube cryocooler.
Demonstrated: 6.6~MV/m at 100\% duty cycle, cryogen-free operation.

\begin{finding}[Nb$_3$Sn Phase~III Fabrication Update]
\label{find:nb3sn_update}
\confirmed{Co-sputtering reduces Nb$_3$Sn fabrication risk for
Phase~III. Combined with compact cryomodule development, the 256-cavity
array at 4~K is now on a credible engineering pathway. Phase~III
cost estimate remains \$75--230M savings over Nb at 20~mK (W4i13),
with reduced fabrication risk.}

\newfinding{Cu-substrate-compatible Nb$_3$Sn co-sputtering could
enable lower-cost cavity bodies for Phase~III. Nb cavities cost
\$50--100K each; Cu cavities with Nb$_3$Sn coating could cost
\$20--40K each. At 256 cavities: potential additional saving of
\$8--15M on cavity bodies.}
\end{finding}


%=============================================================================
\section{Finding 4: Interface Piezoelectricity at Al/Si Junctions}
\label{sec:interface_piezo}
%=============================================================================

\subsection{Discovery and Quantification}

Jahanbani et al.\ (Nature Communications, December 2025) used
aluminum interdigital transducers (IDTs) on high-resistivity silicon
to demonstrate piezoelectric transduction from room temperature to
millikelvin temperatures. The effective coupling:
\begin{equation}
K^2_{\rm eff} = (3 \pm 0.4) \times 10^{-5}\%
= 3\ee{-7},
\label{eq:k2_interface}
\end{equation}
which is comparable to weakly piezoelectric substrates like quartz
($K^2 \sim 10^{-4}$) but orders of magnitude below strong
piezoelectrics like LiNbO$_3$ ($K^2 \sim 5\%$) or AlN ($K^2 \sim 7\%$).

The physical origin is the broken centrosymmetry at the Al/Si
interface. Although bulk Si is centrosymmetric (diamond cubic,
$O_h$ point group, no piezoelectricity), the interface reduces
symmetry to $C_{4v}$, permitting a nonzero piezoelectric tensor.

\subsection{Impact on Crystalline Si PnC Resonators}

For a soft-clamped PnC defect mode with surface participation ratio
$p_{\rm surf}$, the piezoelectric quality factor limit is:
\begin{equation}
Q_{\rm piezo}^{-1} \approx K^2_{\rm eff} \cdot p_{\rm surf}
\cdot \frac{\rho_{\rm EM}(\omega_m)}{1},
\label{eq:q_piezo}
\end{equation}
where $\rho_{\rm EM}(\omega_m)$ is the electromagnetic mode density
at the mechanical frequency. In a well-designed PnC, $p_{\rm surf}$
is exponentially suppressed ($\sim 10^{-4}$ to $10^{-6}$). Taking
$p_{\rm surf} = 10^{-4}$ (conservative for soft clamping):
\begin{equation}
Q_{\rm piezo} \sim \frac{1}{3\ee{-7} \times 10^{-4}}
= 3.3\ee{10}.
\label{eq:q_piezo_num}
\end{equation}

This is close to (but above) the demonstrated $Q > 10^{10}$ of
crystalline Si PnC resonators. With better soft clamping
($p_{\rm surf} \sim 10^{-6}$): $Q_{\rm piezo} \sim 3\ee{12}$,
which is irrelevant.

\subsection{Revised SNR Estimate}

W4i13 claimed a combined $430\times$ SNR improvement for cryogenic
crystalline Si over room-temperature Si$_3$N$_4$ (from $10\times Q$
improvement and $43\times$ thermal noise reduction). If interface
piezoelectricity limits $Q$ to $3\ee{10}$ (with Al metallisation
and conservative $p_{\rm surf}$):

\begin{equation}
\frac{\SNR_{\rm cryo\;Si}}{\SNR_{\rm RT\;Si_3N_4}}
= \frac{Q_{\rm Si,\;limited}}{Q_{\rm Si_3N_4}} \times
\frac{\bar{n}_{\rm th}(300~{\rm K})}{\bar{n}_{\rm th}(7~{\rm K})}
= \frac{3\ee{10}}{10^9} \times 43
= 1290.
\label{eq:snr_revised}
\end{equation}

This is actually \textit{larger} than the W4i13 estimate of $430\times$
because the piezoelectric $Q$ limit ($3\ee{10}$) exceeds the
demonstrated $Q$ ($10^{10}$).

\begin{finding}[Interface Piezoelectricity Impact]
\label{find:interface_piezo}
\confirmed{Interface piezoelectricity at Al/Si junctions does not
degrade the DFD crystalline Si MEMS programme. Even in the most
conservative scenario ($Q$ limited to $3\ee{10}$ by Al metallisation),
the $Q/\Qpsi$ margin remains $136\times$ and the SNR improvement
over room-temperature Si$_3$N$_4$ is $\sim 1300\times$.}

\caution{Mitigation is still recommended: use optical (non-metallic)
readout for the highest-$Q$ measurements, or ensure no Al metallisation
contacts the PnC defect mode region. This is standard practice in
high-$Q$ nanomechanics.}

\newfinding{The same interface piezoelectricity provides an additional
transduction mechanism for MEMS readout. While $K^2 = 3\ee{-7}$ is
too weak for efficient transduction, it provides a diagnostic tool
for confirming mechanical mode frequencies without optical access.}
\end{finding}


%=============================================================================
\section{Finding 5: 3C-SiC Multimode Electromechanical Systems}
\label{sec:3c_sic}
%=============================================================================

\subsection{Nature Communications Results (January 2025)}

Li et al.\ demonstrate cubic silicon-carbide (3C-SiC) membrane
crystals with:

\begin{itemize}[nosep]
\item Heteroepitaxial growth on Si substrates produces non-uniform
  built-in tensile stress (from $\sim 20\%$ lattice mismatch after
  partial relaxation)
\item 21 distinct high-$Q$ mechanical modes controlled from a single
  square membrane
\item Degeneracy-breaking: anisotropic stress and crystal symmetry
  break the $C_{4v}$ degeneracy of membrane modes, producing
  resolvable mode pairs/clusters
\item Direct microwave electromechanical coupling to superconducting
  circuits
\item Group delays up to $\sim 1$~hour, indicating extremely low
  dephasing
\end{itemize}

\subsection{Follow-Up: Frequency Stability (npj QI, 2026)}

The follow-up paper demonstrates superior frequency stability and
long-lived state-swapping in 3C-SiC mechanical mode pairs, confirming
the material's potential for quantum information applications.

\subsection{DFD Relevance Assessment}

\textbf{Advantages for DFD}:
\begin{enumerate}[nosep]
\item \textit{Built-in tensile stress}: Unlike monolithic 4H-SiC,
  3C-SiC on Si has stress enabling dissipation dilution.
\item \textit{Multimode operation}: 21 modes from a single membrane
  enables broadband psi-field spectroscopy. The psi-field couples to
  all mechanical modes proportionally to their effective mass and
  frequency; measuring $\Qpsi$ across multiple modes provides
  frequency-resolved data for the in-gap diagnostic.
\item \textit{Native electromechanical coupling}: Direct coupling to
  superconducting circuits eliminates the optical readout chain,
  simplifying the Phase~II/III experimental setup.
\item \textit{Radiation hardness}: 3C-SiC's wide bandgap and chemical
  stability are advantageous for long-term operation in radiation
  environments (relevant for ESS Channel B, P6 space applications).
\end{enumerate}

\textbf{Limitations}:
\begin{enumerate}[nosep]
\item Individual mode $Q$ values have not been characterised at the
  $10^9$--$10^{10}$ level; the ``high-$Q$'' characterisation in
  the Nature Communications paper refers to $Q > 10^5$--$10^6$.
\item Defect density in 3C-SiC is higher than in 4H-SiC due to
  stacking faults from heteroepitaxial growth.
\item Soft-clamped PnC geometry has not been applied to 3C-SiC yet.
\end{enumerate}

\begin{finding}[3C-SiC Multimode Assessment]
\label{find:3c_sic}
\verdict{3C-SiC membrane crystals are a genuinely new option for
DFD multimode psi-detection. Add to materials watch list at MEDIUM
priority. The multimode + electromechanical coupling combination is
unique among all assessed materials. Key milestone to watch:
demonstration of soft-clamped PnC geometry in 3C-SiC with $Q > 10^8$
per mode.}
\end{finding}


%=============================================================================
\section{Additional Assessments}
\label{sec:additional}
%=============================================================================

\subsection{2D Material Resonators (Graphene, MoS$_2$) --- KILLED}
\label{sec:2d_killed}

Recent progress in 2D NEMS resonators:
\begin{itemize}[nosep]
\item \textbf{MoS$_2$}: Strain dilution of thermoelastic damping
  raises $Q$ from 55 to 3211 at room temperature (Nano Letters 2026,
  26(9), 3141--3148). Gate-induced tensile strain provides $58\times$
  enhancement.
\item \textbf{Graphene}: $Q \sim 10^6$ at 15~mK (ground-state cooled);
  $Q \sim 10^4$--$10^5$ at 4~K.
\item \textbf{WSe$_2$}: $Q = 4.7\ee{4}$ at 4~K (monolayer).
\item \textbf{MoS$_2$ arrays}: Mass transfer printing enables
  scalable NEMS arrays (Microsystems \& Nanoengineering, January 2026).
\end{itemize}

\textbf{DFD assessment}: The best 2D resonator $Q$ values are:
\begin{center}
\begin{tabular}{lcc}
\toprule
Material & Best $Q$ & Temperature \\
\midrule
MoS$_2$ & 3,211 & 300~K \\
Graphene & $\sim 10^6$ & 15~mK \\
WSe$_2$ & $4.7\ee{4}$ & 4~K \\
\bottomrule
\end{tabular}
\end{center}

Even the best value ($Q \sim 10^6$ for graphene at 15~mK) is
$1000\times$ below Si$_3$N$_4$ at room temperature and $10^4\times$
below crystalline Si at 7~K. The $\Qpsi$ margin at $Q = 10^6$ is
$10^6/2.2\ee{8} = 0.005$ --- \textbf{below the detection threshold
by $200\times$}.

\begin{finding}[2D Material Resonators]
\label{find:2d_killed}
\killed{All 2D material resonators (graphene, MoS$_2$, WSe$_2$, etc.)
for DFD psi-detection. Best $Q \sim 10^6$ (graphene, 15~mK) is
$200\times$ below the $\Qpsi$ threshold and $10^4\times$ below
crystalline Si. The fundamental limitation is that atomically thin
membranes have extremely large surface-to-volume ratios, making them
dominated by surface losses (adsorbates, defects, polymer residue
from transfer). No plausible path to $Q > 10^8$ exists for any 2D
material.}
\end{finding}

\subsection{High-$T_c$ Superconductors for SRF (MgB$_2$, YBCO/REBCO)}
\label{sec:high_tc}

\textbf{MgB$_2$} ($T_c = 39$~K): Interest for SRF as it could
operate at 20~K on cryocoolers. Current status: surface resistance
$R_s \sim 1$--$10$~$\mu\Omega$ at GHz frequencies, compared to
$R_s \sim 10$~n$\Omega$ for Nb at 2~K. The $Q_0$ gap is
$\sim 100$--$1000\times$ below Nb. Active research at Frontiers
topic ``Progress on Superconducting Materials for SRF Applications''
(2023--2025), but no breakthrough approaching Nb performance.

\textbf{YBCO/REBCO} ($T_c \sim 92$~K): Could operate at 50--77~K.
Microwave surface resistance measured via Hakki-Coleman resonators at
$\sim 8$~GHz. Demonstrated: $R_s \sim 0.1$--$1$~m$\Omega$ at 77~K,
$\sim 10^4\times$ worse than Nb at 2~K. The lossy grain boundaries
and anisotropic gap of $d$-wave cuprates make high-$Q$ SRF cavities
fundamentally more difficult than for $s$-wave superconductors
(Nb, Nb$_3$Sn, MgB$_2$).

W4i10 already noted that YBCO London penetration depth limitation
makes it \textit{worse} than Nb for psi-coupling (YBCO 100$\times$
enhancement DEMOTED to $<0.001\times$).

\begin{finding}[High-$T_c$ SRF Assessment]
\label{find:high_tc}
\killed{High-$T_c$ superconductors (MgB$_2$, YBCO, REBCO) for DFD
SRF cavities. Surface resistance is $100$--$10^4\times$ worse than Nb
at their respective operating temperatures. The $d$-wave gap symmetry
of cuprates fundamentally limits microwave $Q$. MgB$_2$ ($s$-wave)
is the most promising but still $\sim 100\times$ worse than Nb at
equivalent BCS conditions. Nb$_3$Sn at 4~K remains the optimal
high-temperature SRF material for DFD Phase~III.}

\honest{Room-temperature-adjacent SRF is a physics impossibility with
current materials. The BCS surface resistance has an exponential
temperature dependence: $R_{\rm BCS} \propto \exp(-\Delta/k_BT)$.
Even a hypothetical $T_c = 300$~K superconductor with gap ratio
$\Delta/k_BT_c = 1.76$ would have $R_{\rm BCS}(300~{\rm K}) \propto
\exp(-1.76) \sim 0.17$, giving finite surface resistance (not
negligible). Room-temperature SRF at $Q > 10^{10}$ requires
$R_s < 10$~n$\Omega$, which demands $T/T_c \ll 1$ regardless of
$T_c$ value.}
\end{finding}

\subsection{AlN Piezo Readout --- Expanded Assessment}
\label{sec:aln_update}

Two new developments relevant to AlN readout for DFD:

\begin{enumerate}[nosep]
\item \textbf{Scandium-doped AlN (ScAlN)}: Sc doping increases the
  piezoelectric coefficient of AlN by up to $5\times$
  ($d_{33} \sim 25$~pm/V for Sc$_{0.36}$Al$_{0.64}$N vs
  $\sim 5$~pm/V for pure AlN). Recent review (Microsystems \&
  Nanoengineering, 2025) confirms ScAlN MEMS mirrors with enhanced
  piezo response. For DFD: ScAlN PnC membranes could improve the
  piezo readout sensitivity by $5\times$ over pure AlN, narrowing
  the gap with optical interferometry.

\item \textbf{All-nitride superconducting qubits (ALDmons)}: AlN
  is being used as a tunnel barrier in transmon qubits, with atomic
  layer deposition enabling critical current density tuning over 7
  orders of magnitude (Quantum Zeitgeist, 2025). The piezoelectric
  property of AlN in Josephson junctions is a \textit{loss source}
  for qubits but a \textit{feature} for DFD --- direct electromechanical
  coupling between the qubit and a mechanical mode.
\end{enumerate}

\begin{finding}[AlN Readout Update]
\label{find:aln_update}
\verdict{AlN/ScAlN piezoelectric readout confirmed as a viable
Phase~II option. ScAlN provides $5\times$ enhanced coupling. The
emerging all-nitride qubit platform creates a natural integration
pathway for AlN MEMS + superconducting readout. No change to primary
material selection (Si$_3$N$_4$ sensor + AlN/ScAlN readout is the
recommended hybrid architecture for Phase~II).}
\end{finding}


%=============================================================================
\section{Updated Materials Roadmap}
\label{sec:roadmap}
%=============================================================================

\begin{table}[H]
\centering
\caption{Updated materials roadmap (W4i14). Changes from W4i13 in
\textcolor{teal}{teal}.}
\label{tab:roadmap}
\renewcommand{\arraystretch}{1.3}
\begin{tabular}{p{2cm}p{3cm}p{2.5cm}p{2cm}p{3cm}}
\toprule
\textbf{Phase} & \textbf{SRF Material} & \textbf{MEMS Material} &
\textbf{$T_{\rm op}$} & \textbf{Change from W4i13} \\
\midrule
I (2027--28) &
  N-doped Nb + \textcolor{teal}{O-doped Nb (co-primary)} &
  Si$_3$N$_4$ PnC (room $T$) &
  SRF: 20~mK; MEMS: 300~K &
  \textcolor{teal}{O-doped upgraded to co-primary; mixed
  N/O 4-cavity array recommended} \\
\midrule
II (2028--29) &
  N-doped Nb &
  Si$_3$N$_4$ PnC + crystalline Si PnC (cryo);
  \textcolor{teal}{AlN/ScAlN hybrid readout} &
  SRF: 20~mK; MEMS: 300~K / 7~K &
  \textcolor{teal}{ScAlN readout pathway added} \\
\midrule
III (2030+) &
  Nb$_3$Sn \textcolor{teal}{(co-sputtered)} &
  Best of: cm-scale Si$_3$N$_4$, cryo Si,
  \textcolor{teal}{3C-SiC multimode (if $Q > 10^8$ demonstrated)} &
  SRF: 4~K; MEMS: TBD &
  \textcolor{teal}{4H-SiC removed; 3C-SiC added; co-sputtering
  specified for Nb$_3$Sn; Cu substrate option noted} \\
\bottomrule
\end{tabular}
\end{table}


%=============================================================================
\section{Consistency Checks and Flagged Issues}
\label{sec:consistency}
%=============================================================================

\subsection{Resolved: W4i13 4H-SiC Quality Factor Discrepancy}

W4i13 cited $Q_{\rm int} = 2\ee{7}$ for 4H-SiC (attributed to
Sementilli et al., Nano Letters 2025). The Phys.\ Rev.\ Applied
paper (September 2025) by the same group reports $Q_{\rm int}$ up to
$2\ee{5}$ at room temperature for nanostrings. The discrepancy may
arise from: (a) different device geometries (nanostring vs trampoline),
(b) different frequency regimes (the Nano Letters paper may report
$Q$ at lower frequencies where thermoelastic damping is less
important), or (c) different definitions of ``intrinsic'' $Q$.

\caution{The most conservative and well-characterised value is
$Q_{\rm int} = 2\ee{5}$ from the Phys.\ Rev.\ Applied paper, which
is a comprehensive study with SLDV characterisation. Previous claims of
$Q = 2\ee{7}$ may reflect measurements at specific favorable frequencies
or geometries. For DFD programme planning, use $Q_{\rm int} = 2\ee{5}$
as the baseline for unstressed 4H-SiC.}

\subsection{Unresolved: Phase~I $Q_0$ Sensitivity Scaling (Inherited from W4i13)}

W4i13 flagged an inconsistency: does $\lam_{\rm min} \propto 1/Q_0$
or $\propto 1/\sqrt{Q_0}$? This depends on whether Phase~I operates
in the $C \ll 1$ or $C \sim 1$ regime. At $Q_0 = 5\ee{10}$ (baseline):
\begin{equation}
C = \frac{\lam^2 Q_0^2 F_\psi}{\Meff \omega^2 / c^4}.
\end{equation}

From MASTER\_FINDINGS\_CORPUS: at the SQMS bound $\lam = 1.56\ee{-9}$
and $Q_0 = 5\ee{10}$, the cooperativity $C$ is expected to be $\ll 1$
(otherwise detection would already have occurred). Therefore,
$\lam_{\rm min} \propto 1/Q_0$ is the correct scaling for Phase~I.

At $Q_0 = 2\ee{11}$ (target), $C$ increases by $(2\ee{11}/5\ee{10})^2
= 16\times$. If baseline $C \sim 10^{-3}$, then target $C \sim 0.016$
--- still $C \ll 1$. The $1/Q_0$ scaling holds throughout Phase~I.

\confirmed{Phase~I operates in the $C \ll 1$ regime at both baseline
and target $Q_0$. The sensitivity scaling $\lam_{\rm min} \propto
1/Q_0$ from W4i10 Eq.~(5) is correct. The W4i13 flag is resolved.}

\subsection{New Flag: Crystalline Si PnC Fabrication Availability}

W4i13 recommended contacting EPFL (Kippenberg/Ghadimi) or DTU Nanolab
for crystalline Si PnC devices. The original Ghadimi et al.\ result
(Nature Physics, 2022) used a specific fabrication protocol: SOI
(silicon-on-insulator) wafer with 220~nm device layer, e-beam
lithography for PnC pattern, RIE dry etch, HF vapor release of
buried oxide, and tensile stress from differential thermal contraction.

\caution{This fabrication requires:
\begin{enumerate}[nosep]
\item SOI wafers with specific device layer thickness (220~nm) and
  stress engineering
\item E-beam lithography with $\sim 50$~nm feature resolution for
  PnC bandgap optimisation
\item Critical point drying or HF vapor release to avoid stiction
\item Cryogenic characterisation at $< 10$~K
\end{enumerate}
All steps are standard in academic nanofabrication but require
specialised equipment. The DFD programme should budget 6--12 months
and \$100--200K for first device procurement, consistent with W4i13
recommendation. Key risk: stress control in the released Si membrane
may not match the Ghadimi recipe exactly, requiring optimisation
iterations.}


%=============================================================================
\section{Summary of Material Status Changes}
\label{sec:changes}
%=============================================================================

\begin{table}[H]
\centering
\caption{Material status changes from W4i13 to W4i14.}
\renewcommand{\arraystretch}{1.3}
\begin{tabular}{p{4cm}p{2.5cm}p{2.5cm}p{4cm}}
\toprule
\textbf{Material/Technology} & \textbf{W4i13 Status} &
\textbf{W4i14 Status} & \textbf{Reason} \\
\midrule
4H-SiC (monolithic) & WATCH: Long-term & \textbf{DOWNGRADED:
Speculative} & No built-in stress; dissipation dilution blocked;
$Q_{\rm int} = 2\ee{5}$ not $2\ee{7}$ \\
\midrule
3C-SiC membrane & Not assessed & \textbf{ADDED: WATCH (Medium)} &
Built-in stress; 21 multimode electromech coupling; direct SC
circuit interface \\
\midrule
O-doped Nb & Phase~I backup & \textbf{UPGRADED: Phase~I co-primary} &
APL March 2026 quantitative comparison; superior high-field stability \\
\midrule
Nb$_3$Sn (co-sputtered) & Phase~III (vapor diffusion assumed) &
\textbf{UPGRADED: Co-sputtering preferred} & Improved homogeneity,
reproducibility; Cu substrate option \\
\midrule
Interface piezoelectricity & Not assessed & \textbf{FLAGGED:
Design constraint} & Al/Si coupling $K^2 = 3\ee{-7}$ limits $Q$ to
$\sim 3\ee{10}$ if Al used; mitigable by optical readout \\
\midrule
ScAlN readout & Not assessed & \textbf{ADDED: Phase~II option} &
$5\times$ piezo enhancement over AlN; integration with all-nitride
qubit platform \\
\midrule
2D resonators & Not assessed & \textbf{KILLED} & Best $Q \sim 10^6$
(graphene, 15~mK); $200\times$ below $\Qpsi$ threshold \\
\midrule
High-$T_c$ SRF & Not assessed & \textbf{KILLED} & MgB$_2$: $100\times$
worse $R_s$; YBCO: $10^4\times$ worse; room-temp SRF impossible \\
\midrule
Phase~I $Q_0$ scaling & Flagged (W4i13) & \textbf{RESOLVED} &
$C \ll 1$ at both baseline and target $Q_0$; $1/Q_0$ scaling correct \\
\bottomrule
\end{tabular}
\end{table}


%=============================================================================
\section{Recommendations for Next Iteration}
\label{sec:next}
%=============================================================================

\begin{enumerate}[nosep]
\item \textbf{Crystalline Si MEMS procurement} (inherited from W4i13,
  unchanged): Contact EPFL (Kippenberg) or DTU Nanolab. Budget:
  \$100--200K, 6--12 months to first device.

\item \textbf{O-doped cavity request}: Request SQMS include O-doped
  Nb cavities in the Phase~I 4-cavity array alongside N-doped.
  Cost: \$0 within programme. Provides doping cross-check on Q-ratio
  diagnostic.

\item \textbf{Nb$_3$Sn co-sputtering CDR}: Commission a conceptual
  design report for 256-cavity co-sputtered Nb$_3$Sn array at 4~K.
  Include Cu-substrate cost analysis. Engage Fermilab SRF/SQMS.
  Cost: \$50--100K.

\item \textbf{3C-SiC collaboration}: Identify groups fabricating
  3C-SiC membrane crystals (Aalto, NTT, ETH) for potential Phase~III
  multimode psi-spectroscopy. No procurement needed yet; literature
  monitoring and initial discussion.

\item \textbf{Interface piezoelectricity mitigation}: Ensure
  crystalline Si PnC resonator design avoids Al metallisation in
  the defect mode region. Optical readout is already the baseline;
  confirm this is maintained for cryogenic measurements.

\item \textbf{Standard territory exhausted}: The materials science
  survey is approaching convergence. The material hierarchy is now:
  \begin{center}
  Si$_3$N$_4$ (Phase II, proven) $>$ Crystalline Si (Phase II+, 7K)
  $>$ Nb$_3$Sn (Phase III SRF) $>$ 3C-SiC (Phase III multimode,
  speculative) $\gg$ [everything else KILLED]
  \end{center}
  Further iterations should focus on \textit{engineering} (fabrication
  protocols, integration, procurement) rather than \textit{survey}
  (new materials). The search space is exhausted.
\end{enumerate}


\end{document}
